\section{Phase 1 results: accuracy on the GPU}\label{sec:gpu}

\subsection{What was run}

\file{experiments/run.py} runs four stages for each dataset on an NVIDIA L40 GPU:
\begin{enumerate}
  \item \textbf{errors}: the multiplier error statistics of Section~\ref{sec:errors}.
  \item \textbf{train}: normal float32 training of both networks (MNIST: 15 epochs, Fashion-MNIST:
        30 epochs; Adam, learning rate $10^{-3}$ with cosine decay, batch 128, seed 0).
  \item \textbf{eval}: every trained network in four number formats (float, trunc, int8, int16) times eight
        multipliers, on all 10{,}000 test pictures.
  \item \textbf{finetune}: for int8 and int16, a few more epochs with each multiplier, then evaluate again.
\end{enumerate}
The results are in \file{results/mnist/} and \file{results/fashion/} (\file{eval\_*.csv},
\file{finetune\_*.csv}, \file{multiplier\_errors.csv}).

\subsection{Number formats: scaling is essential}

\begin{figure}[H]
\centering
\includegraphics[width=0.85\linewidth]{domains}
\caption{Test accuracy with the exact multiplier in four number formats. Cutting decimals to
integers without a scale (``cut to int'') collapses to 10\%, which is guessing. int8 and int16 are as good
as float32.}
\label{fig:domains}
\end{figure}

\subsection{Accuracy with every multiplier}

\begin{figure}[H]
\centering
\includegraphics[width=\linewidth]{gpu_acc}
\caption{Test accuracy (10{,}000 test pictures) for int8 networks with each multiplier, before (solid) and
after (light) fine-tuning. The dashed line is the original float32 network. \textbf{Look at the
vertical axis}: each panel spans only about one percentage point, so the bars look more different than they
are.}
\label{fig:gpuacc}
\end{figure}

\begin{table}[H]
\centering
\small
\begin{tabular}{lcccccccc}
\toprule
 & \multicolumn{2}{c}{MNIST MLP} & \multicolumn{2}{c}{MNIST LeNet-5} & \multicolumn{2}{c}{Fashion MLP} &
   \multicolumn{2}{c}{Fashion LeNet-5}\\
\cmidrule(lr){2-3}\cmidrule(lr){4-5}\cmidrule(lr){6-7}\cmidrule(lr){8-9}
 & int8 & +FT & int8 & +FT & int8 & +FT & int8 & +FT\\
\midrule
float32 & \multicolumn{2}{c}{98.42} & \multicolumn{2}{c}{99.20} & \multicolumn{2}{c}{89.79} & \multicolumn{2}{c}{90.99}\\
\midrule
Exact & 98.43 & 98.43 & 99.18 & 99.21 & 89.82 & 89.77 & 90.95 & 90.88\\
Mitchell & 98.46 & 98.46 & 99.12 & 99.20 & 89.75 & 89.82 & 90.82 & 91.15\\
DRUM6 & 98.43 & 98.44 & 99.17 & 99.20 & 89.69 & 89.77 & 90.92 & 90.97\\
RoBA-Dn1 & 98.42 & 98.44 & 99.19 & 99.18 & 89.69 & 89.66 & \textbf{90.20} & 90.77\\
RoBA-UpDn & 98.44 & 98.42 & 99.20 & 99.22 & 89.89 & 89.74 & 90.72 & 90.84\\
RoBA-Dn2 & 98.44 & 98.51 & 99.20 & 99.21 & 89.74 & 89.71 & 90.91 & 91.01\\
RoBA-UpDnDn & 98.43 & 98.44 & 99.19 & 99.18 & 89.83 & 89.75 & 91.07 & 90.95\\
RoBA-(UpDn)2 & 98.44 & 98.47 & 99.21 & 99.19 & 89.77 & 89.73 & 90.98 & 90.97\\
\bottomrule
\end{tabular}
\caption{Test accuracy in \% (int8). ``+FT'' = after fine-tuning with that multiplier. The biggest
drop anywhere is RoBA-Dn1 on Fashion-MNIST LeNet-5: $90.95\to90.20$, i.e.\ $-0.75$ points, and
fine-tuning brings it back to 90.77.}
\label{tab:gpuacc}
\end{table}

\begin{careful}
We trained with one random seed. Training the same network twice with different seeds usually changes
accuracy by about 0.1 points, so differences smaller than that (most of this table) are noise, not a real
effect. ``Mitchell + FT is the best on Fashion LeNet-5'' is \emph{not} a real finding.
\end{careful}

\subsection{Why can a 9\%-wrong multiplier keep 99\% of the accuracy?}\label{sec:whyok}

RoBA-Dn1 is on average 8.9\% too small on every single multiplication, yet the networks barely notice.
Three reasons:

\paragraph{1. A consistent ``too small'' is like turning the volume down.} If every product in a sum is
about 9\% too small, the whole sum is about 9\% too small. The \emph{largest} score stays the largest.
\begin{example}[: shrinking all scores]
Scores $(5.0, 4.0, 1.0)$, answer: class 0. Multiply all by $0.91$: $(4.55, 3.64, 0.91)$, answer: still
class 0. It only goes wrong when two scores are nearly tied \emph{and} the bias (which is not shrunk)
tips the balance, which is rare. This is also why the loss of Dn1 is \emph{lower} than exact on Fashion
(0.376 vs 0.434 for the MLP): smaller scores make the network less over-confident when it is wrong.
\end{example}

\paragraph{2. Random errors cancel in a long sum.} A neuron in the MLP's first layer adds 784 products.
Errors that are sometimes positive and sometimes negative (UpDn, DRUM6) largely cancel out.
\begin{example}[: errors cancelling]
Four products whose true value is 100 each come out as $106, 96, 103, 95$ (errors $+6\%, -4\%,
+3\%, -5\%$). Their sum is $106+96+103+95 = 400$: exactly the true sum. Each one is up to 6\% wrong,
but together they are 0\% wrong. In general, the
relative error of a sum of $n$ independent errors shrinks about like $1/\sqrt{n}$: for $n=784$, about 28
times smaller.
\end{example}

\paragraph{3. Networks are robust by nature.} They are trained to recognise a shoe whether it is
slightly brighter, shifted or noisy. A small error in the arithmetic looks to them like a bit of noise.

\begin{keyidea}
Accuracy is \textbf{not} the problem: every multiplier stays within 0.75 points of exact, and most within
0.2 points. So the choice of multiplier should be made on \textbf{speed and energy}, which is what Phase~2
measures on the real CPU.
\end{keyidea}

\subsection{What the old cost estimate said}

Before measuring on Ibex, \file{roba/cost.py} estimated cycles by reading the cycles per multiply from
Fig.~4 of the paper (e.g.\ 8-bit: exact 44, Dn1 19) and multiplying by the number of MACs. Those
estimates are kept in the CSV files (\code{mult\_cycles}, \code{energy\_rel}) but are replaced by the
real measurements below. The measurements turned out higher than the estimates, because a real program
also spends time loading numbers, handling signs and looping, not only multiplying.
\newpage
